% Source-aligned English for printed p.175.
\SourcePage{175}

\LayerTarget{tohoku:construction:cartan-leray-fixed-cover:p2}
\noindent $R(H^0C)$ of the functor $H^0C$, whose initial term is
\LayerTarget{tohoku:eq:cartan-leray-abstract-E2}
$E_2^{pq}=H^p(R^qC)$.
\LayerTarget{tohoku:proof:cartan-leray-derived-functor}
This may be seen directly by taking an injective resolution of $F\in\mathbf C$ and
examining the spectral sequences of the double complex obtained by applying
$\mathbf C$ to that resolution; or better, by observing that $H^0(K)$, regarded as a
covariant left-exact functor on $\mathbf C'$ (with values in the category of abelian
groups), has right derived functors $H^p(K)$. Thus our spectral sequence is a
simple special case of Theorem 2.4.1. Of course, $(R^qC)(F)$ is the complex
whose components are $(R^qC^p)(F)$, if the $C^p$ are the components of $C$.
It remains, in the present case, to describe the $R^qC^p$. By
\LayerTarget{tohoku:eq:cover-cochains-product}
\[
C^pF=\prod_{\sigma^p}\Gamma(U_{\sigma^p},F)
\]
(the product extending over all sequences $\sigma^p=(i_0,\ldots\mathpunct{.},i_p)$ of
$p+1$ indices of the covering $U=(U_i)_{i\in I}$), one immediately
sees that
\LayerTarget{tohoku:eq:derived-cover-cochains-product}
\[
R^qC^p(F)=\prod_{\sigma^p}R^q(\Gamma(U_{\sigma^p},F)).
\]
\LayerTarget{tohoku:lemma:open-sections-derived}
Now if $V$ is an open subset of $X$, the right derived functors of
$\Gamma(V,F)=\Gamma(F\mid V)$ are easily made explicit, since the restriction
functor $F\to F\mid V$ from $\mathbf C^X$ to $\mathbf C^V$ is exact and sends
injective objects to injective objects (Proposition 3.1.3). Thus
\LayerTarget{tohoku:eq:open-sections-derived}
$R^q(\Gamma(V,F))=(R^q\Gamma^V)(F\mid V)=H^q(V,F)$.
\LayerTarget{tohoku:def:cohomology-presheaf}
Let $H^q(F)$ denote the \textit{presheaf} on $X$ whose value on an open subset
$V$ is $H^q(V,F)$. Then
\LayerTarget{tohoku:eq:derived-cover-cochains-presheaf}
\[
R^qC^p(F)=\prod_{\sigma^p}\Gamma(U_{\sigma^p},H^q(F))
=C^p(\mathbf U,H^q(F)).
\]
Of course, the differential on $R^qC(F)=\displaystyle\sum\limits_p R^qC^p(F)$ is that of
$C(\mathbf U,H^q(F))$. Consequently,
\LayerTarget{tohoku:eq:cartan-leray-E2}
$E_2^{p,q}(F)=H^p(\mathbf U,H^q(F))$.
As for the abutment of the spectral sequence, it is the right derived functor
of $H^0C(F)=\Gamma(X,F)$, namely $(H^n(X,F))$. Hence:

\LayerTarget{tohoku:theorem:3.8.1}
\textsc{Theorem 3.8.1.} \textit{Let $X$ be a topological space
endowed with an open covering $\mathbf U$. Then there exists a cohomological
spectral functor on the category $\mathbf C^X$ of sheaves of abelian groups on $X$,
abutting to the graded functor $(H^n(X,F))$, whose initial term is}
\LayerTarget{tohoku:eq:3.8.1-E2}
\[
E_2^{p,q}=H^p(\mathbf U,H^q(F))
\]
\textit{where, for every sheaf $F\in\mathbf C^X$, $H^q(F)$ denotes the presheaf
$V\to H^q(V,F)$ on $X$.}

\LayerTarget{tohoku:remark:3.8.1-no-hypotheses}
Note that this spectral sequence is established here without any paracompactness
hypothesis on $X$ or local-finiteness hypothesis on $\mathbf U$. The
preceding spectral sequence gives functorial homomorphisms
\LayerTarget{tohoku:map:cech-to-derived-cohomology}
\[
H^p(\mathbf U,F)\to H^p(X,F)
\]
and moreover:

\LayerTarget{tohoku:corollary:3.8.1:1:p1}
\textsc{Corollary 1.} \textit{The preceding homomorphisms are
isomorphisms if}

\TohokuContinuation{176}
